\documentclass[12pt]{article}
\usepackage[a4paper,left=25mm,right=25mm,top=15mm,bottom=25mm]{geometry}
\usepackage{graphicx,url}
\usepackage{color}
\usepackage{amssymb,listings}
\usepackage[utf8]{inputenc}
\usepackage[brazilian]{babel}
\usepackage[T1]{fontenc}

% Setting configuration for the text format
%\renewcommand{\contentsname}{Table of Contents}
%\renewcommand{\bibname}{References}
%\titleformat{\chapter}[display]{\normalfont\huge\bfseries}{\filleft\chaptername\ \thechapter}{5pt}{\filleft\Huge}
%\sloppy

\title{Algoritmos e Linguagens de Programação - Reposição Prática}
\author{Rogério Eduardo da Silva \& Claudio Cesar de Sá}
\date{\today}

\graphicspath{{/img/}}   
\DeclareGraphicsExtensions{{.jpg},{.png}}

\lstset {
  numbers=left,
}

\begin{document}
%\pagestyle{empty}
\maketitle



\begin{flushleft}
\textbf{Nome}: \rule{10cm}{0.3mm} 
\end{flushleft}
%\textbf{Matrícula}: \rule{7cm}{0.3mm} 
% %\rule{\linewidth}{1.

\vskip 0.5cm


\begin{enumerate}
%\setlength{\itemsep}{-1pt}

\item Faça um programa em C que, após coletar as respostas acerca dos sintomas de um paciente, auxilie um médico na tomada de decisão da provável 
doença.
\begin{center}
\includegraphics[width=\linewidth]{img/doencas}
\end{center}

{\footnotesize DICA: Uma idéia para visualizar o problema: troque \texttt{sim} por \texttt{1} e \texttt{não} por \texttt{2}. Assim uma entrada igual a: 
\texttt{2 1 1} a resposta terá que ser \texttt{fibromialgia},
\texttt{1 2 1} a resposta terá que ser \texttt{gripe}, e assim sucessivamente.}

\item Usando o \texttt{switch}, faça um programa em C que solicite ao usuário escolher entre três objetos: esfera (e), cilindro (c) ou paralelepípedo (p). 
Em seguida, o programa deve solicitar  as dimensões do objeto, e calcular o volume do mesmo. Sendo:
\begin{description}
\item[Esfera] $Volume = \frac{4}{3} \times \left ( \pi \times r^3 \right )$ -- necessita apenas do raio

\item[Cilindro] $Volume = \pi \times r^2 \times h$
 -- necessita  do raio e altura

\item[Paralelepípedo] $Volume = largura \times altura \times comprimento$  -- necessita de 3 valores

\end{description}

\end{enumerate}
\begin{center}

{\Large Deixe tudo pronto para entradas e saídas via arquivos}\\


\underline{\textbf{Assim que o código estiver pronto chame o professor para apresentação!}}



\end{center}
\end{document}
